La gestion eficiente de inventarios supone un desafio constante en diversos sectores. Problemas como sufrir de debastabecimiento pueden surgir por distintos factores como retradsos en la entrega por parte de los proovedores, variaciones inesperadas de la demanda, o deficiencias en los procesos de planificacion. Estas situaciones, en muchos casos, se encuentran fuera del control de los administradores, lo que incrementa la complejidad para ellos a la hora de tomar desiciones relacionadas con el abastecimiento.

Actualmente, normalmente en muchas empresas se cuenta con un software para gestion de inventarios tipo ERP, con este software se administra la disponibilidad de los productos mediante el uso de diferentes tipos de movimientos de inventario:

Uno de los tipos de movimientos consiste en las entradas de inventario, las cuales responden a la adquisicion de productos atraves de proovedores y son registradas mediante facturas que incluyen informacion como fecha, proovedor y costo. 

Por otro lado, estan las salidas por consumo, las cuales representan la venta, distribución o consumo de los productos, registrando información relacionada con las cantidades entregadas y el destino de estas. Adicionalmente, existen salidas de inventario asociadas a sucesos como devoluciones, productos dañados, vencimientos, pérdidas o ajustes de inventario, se regista la cantidad de productos que deben de ser retirados del inventario junto con su respectivo motivo. Posteriormente, se realizan ajustes de inventario, los cuales pueden ser tanto positivos como negativos, con el fin de corregir diferencias entre las cantidades resgistradas den el sistema y las existencias reales, garantizando asi la consistencia de la informacion.

En general, un buen analisis de estos movimientos resulta fundamental para comprender el comportamiento del consumo de productos y asi mejorar los procesos de planificacion del abastecimiento.